\documentclass[11pt]{article}
\usepackage[margin=1in]{geometry}
\usepackage{amsmath,amssymb,amsthm}
\usepackage{graphicx}
\usepackage{booktabs}
\usepackage{xcolor}
\usepackage{listings}
\usepackage[numbers]{natbib}
\usepackage{hyperref}

\newtheorem{theorem}{Theorem}
\newtheorem{lemma}[theorem]{Lemma}
\theoremstyle{definition}
\newtheorem{definition}[theorem]{Definition}
\theoremstyle{remark}
\newtheorem*{remark}{Remark}

\emergencystretch=2em

\lstdefinelanguage{lean}{
  morekeywords={lemma,set,split_ifs,obtain,constructor,rintro,fin_cases,simpa,positivity,congr,apply,ext,decide,field_simp,ring,div_nonneg,def,theorem,rcases,haveI,subst,nlinarith,linarith,have,refine,intro,induction,with,simp,only,rw,unfold,exact,example,namespace,end,open,fun,Prop,by,noncomputable,section,import},
  sensitive=true,
  morecomment=[l]{--},
  morecomment=[s]{/-}{-/},
}
\lstset{
  language=lean,
  basicstyle=\ttfamily\footnotesize,
  keywordstyle=\color{blue!60!black}\bfseries,
  commentstyle=\color{gray},
  columns=fullflexible,
  keepspaces=true,
  breaklines=true,
  frame=single,
  extendedchars=true,
  literate=
    {ℝ}{{$\mathbb{R}$}}1
    {ℕ}{{$\mathbb{N}$}}1
    {→}{{$\to$}}1
    {↔}{{$\leftrightarrow$}}1
    {∀}{{$\forall$}}1
    {∃}{{$\exists$}}1
    {∈}{{$\in$}}1
    {∑}{{$\sum$}}1
    {∏}{{$\prod$}}1
    {∧}{{$\wedge$}}1
    {≤}{{$\le$}}1
    {≥}{{$\ge$}}1
    {≠}{{$\ne$}}1
    {⨆}{{$\bigsqcup$}}1
    {∉}{{$\notin$}}1
    {←}{{$\leftarrow$}}1
    {ρ}{{$\rho$}}1
    {⟨}{{$\langle$}}1
    {⟩}{{$\rangle$}}1
    {·}{{$\cdot$}}1
    {𝓝}{{$\mathcal{N}$}}1
    {₀}{{$_0$}}1
    {⁻}{{$^{-}$}}1
    {¹}{{$^{1}$}}1
    {↦}{{$\mapsto$}}1
    {⊢}{{$\vdash$}}1
}

\title{Cyclic Industrial Symbiosis:\\ A Lean Proof of the $N$-Partner Product Condition}
\author{Justin Hart, Viridis LLC}
\date{18 November 2026}

\begin{document}
\maketitle

\begin{abstract}
$N\ge2$ firms are arranged in a cycle. Firm $i$ ships a by-product flow $x_i$ to its successor at
unit cost $c_i>0$ and receives its predecessor's flow, worth $v_i>0$ per unit to it. Its net benefit
is $v_ix_{i-1}-c_ix_i$, with indices taken mod $N$. We give a Lean~4 proof, against Mathlib, of the
sealed statement \texttt{P3.CyclicSymbiosis}: a strictly positive flow that gives every firm a
strictly positive net benefit exists if and only if $\prod_i c_i<\prod_i v_i$. Necessity follows by
multiplying the $N$ strict inequalities $c_ix_i<v_ix_{i-1}$ and cancelling $\prod_i x_i$. The
cancellation works because the predecessor map is a bijection of the cycle. For sufficiency we
construct an explicit flow. With $s=(\prod v/\prod c)^{1/N}>1$, every link carries the ratio
$x_i/x_{i-1}=v_i/(c_is)$, and these ratios multiply to one around the cycle. Each firm then keeps the
fraction $1-1/s$ of the value it receives. For $N=2$ the statement reduces to the product condition
$c_Xc_Y<v_Xv_Y$ of the Reciprocity Corridor (Run-128). We include a Lean derivation of the two-partner
feasibility test from the $N$-partner theorem. The result is a statement about a linear,
one-period model. It says nothing about real costs, values, logistics, contracts or how gains are
shared.
\end{abstract}

\section{Introduction}
In industrial symbiosis one firm's residue becomes another firm's input \citep{chertow2000}. The
canonical example is the Kalundborg network, and many exchanges form chains or loops rather than
isolated pairs. The circular economy depends on such exchanges, in which recovered resources
replace virgin ones \citep{geissdoerfer2017}. A decision kernel that proposes exchanges needs a
test for when a loop of exchanges can make \emph{every} participant strictly better off. Without
one, the kernel could propose a loop that one participant would rationally leave.

For two partners the Viridis canon already has this test. The Reciprocity Corridor (Run-128,
\citealp{hart2026rc}) shows that strictly positive flows benefiting both partners exist if and
only if $c_Xc_Y<v_Xv_Y$. This note proves the $N$-partner version for a cycle:
\[
\exists\,x>0\ \ \forall i:\ v_i\,x_{i-1}-c_i\,x_i>0
\quad\Longleftrightarrow\quad
\prod_{i=0}^{N-1}c_i<\prod_{i=0}^{N-1}v_i .
\]
Locally the condition is surprising, globally it is not. A single link can be a loss in isolation
($v_i<c_i$), and the loop can still be feasible as long as the other links make up the shortfall
\emph{multiplicatively}. Project P3 proposed \texttt{P3.CyclicSymbiosis} in session P3-T001
(5~November 2026), and the program owner sealed it under decision D-014. It is the last of the
four P3 statements.

\paragraph{Contributions.}
\begin{itemize}
\item A Lean~4 proof of \texttt{ViridisStatements.P3.CyclicSymbiosis.Statement}
(Section~\ref{sec:proof}). It depends only on the standard axioms \texttt{propext},
\texttt{Classical.choice} and \texttt{Quot.sound}.
\item A human-readable proof with an explicit beneficial flow (Section~\ref{sec:math}).
\item A Lean corollary that derives the two-partner feasibility test of Run-128 from the
$N$-partner theorem (Section~\ref{sec:n2}).
\item A list of what the model leaves open (Section~\ref{sec:limits}).
\end{itemize}

\section{Background and related work}
\paragraph{Industrial symbiosis and the circular economy.}
\citet{chertow2000} defined industrial symbiosis as the exchange of materials, energy, water and
by-products among firms and gave a taxonomy of exchange types. \citet{geissdoerfer2017} review the
circular economy as a sustainability paradigm built on substituting recovered resources for virgin
ones. Earlier P3 sessions gave Lean proofs of three accounting statements.
\texttt{P3.LoopMassBalance} (P3-T002) says that virgin input is $D(1-r)$ under recovery fraction
$r$. \texttt{P3.CascadedRecovery} (P3-T003) says that the virgin share of $n$ stages is
$\prod_i(1-r_i)$. \texttt{P3.RecoveryWorkFloor} (P3-T004) says that fresh supply is at least
$(T-e_{\max}w_r)/e_{\max}$, building on the Useful and Dissipated Work Ledger
\citep{hart2026udwl}. Those statements are about physical accounting. This one is about incentives.

\paragraph{The Reciprocity Corridor.}
Run-128 \citep{hart2026rc} studies two partners $X,Y$ with flows $x,y>0$ and payoffs
$u_X=v_Xy-c_Xx$ and $u_Y=v_Yx-c_Yy$. Both payoffs are positive exactly when
$c_X/v_X<y/x<v_Y/c_Y$, and that corridor is non-empty exactly when $c_Xc_Y<v_Xv_Y$. The cyclic
model keeps Run-128's standing assumptions: payoffs are linear and separable, there is one period,
disagreement utilities are zero, there are no externalities, and coefficients are known and
positive.

\paragraph{Formal tools.}
The proof is written in Lean~4 \citep{moura2021lean4} (toolchain v4.28.0) against Mathlib
\citep{mathlib2020} (commit \texttt{8f9d9cff}).

\section{The model and the statement}
\label{sec:model}
\begin{definition}[Cyclic symbiosis model]
Partners are indexed by $\mathrm{Fin}\,N$. The predecessor is
$\mathrm{prev}(i)=(i+N-1)\bmod N$. For costs $c$, values $v$ and flow $x$, the net benefit of
partner $i$ is $\mathrm{net}_i=v_i\,x_{\mathrm{prev}(i)}-c_i\,x_i$. A flow is \emph{strictly
beneficial} if $x_i>0$ and $\mathrm{net}_i>0$ for every $i$.
\end{definition}

\begin{theorem}[\texttt{P3.CyclicSymbiosis}]\label{thm:main}
Let $N\ge2$ and $c_i,v_i>0$ for all $i$. A strictly beneficial flow exists if and only if
$\prod_ic_i<\prod_iv_i$.
\end{theorem}

\begin{lstlisting}[caption={Definitions and \texttt{Statement} from the sealed module \texttt{ViridisStatements.P3.CyclicSymbiosis} (docstrings omitted).}]
def prev {N : ℕ} (i : Fin N) : Fin N :=
  ⟨(i.val + (N - 1)) % N, Nat.mod_lt _ (Nat.zero_lt_of_lt i.isLt)⟩

def net {N : ℕ} (c v x : Fin N → ℝ) (i : Fin N) : ℝ := v i * x (prev i) - c i * x i

def StrictlyBeneficialFlow {N : ℕ} (c v x : Fin N → ℝ) : Prop :=
  (∀ i, 0 < x i) ∧ ∀ i, 0 < net c v x i

def CycleProductCondition {N : ℕ} (c v : Fin N → ℝ) : Prop := ∏ i, c i < ∏ i, v i

def Statement : Prop :=
  ∀ (N : ℕ), 2 ≤ N → ∀ c v : Fin N → ℝ, (∀ i, 0 < c i) → (∀ i, 0 < v i) →
    ((∃ x : Fin N → ℝ, StrictlyBeneficialFlow c v x) ↔ CycleProductCondition c v)
\end{lstlisting}

\section{Mathematical proof}
\label{sec:math}
\begin{lemma}[The cycle]\label{lem:prev}
$\mathrm{prev}(0)=N-1$ and $\mathrm{prev}(i)=i-1$ for $1\le i<N$. Hence $\mathrm{prev}$ is a
bijection of $\mathrm{Fin}\,N$, and $\prod_ix_{\mathrm{prev}(i)}=\prod_ix_i$ for every $x$.
\end{lemma}
\begin{proof}
For $i=0$, $(N-1)\bmod N=N-1$. For $i\ge1$, $i+N-1=(i-1)+N$, so the remainder is $i-1$. The
explicit form is injective, and an injective self-map of a finite set is a bijection. Reindexing a
product along a bijection does not change it.
\end{proof}

\begin{proof}[Proof of Theorem~\ref{thm:main}, necessity]
Let $x$ be strictly beneficial. Then $0<c_ix_i<v_ix_{\mathrm{prev}(i)}$ for every $i$. The index set
is non-empty, so multiplying these strict inequalities between positive numbers gives
$\prod_ic_i\prod_ix_i<\prod_iv_i\prod_ix_{\mathrm{prev}(i)}=\prod_iv_i\prod_ix_i$ by
Lemma~\ref{lem:prev}. Since $\prod_ix_i>0$, we can cancel it.
\end{proof}

\begin{proof}[Proof of Theorem~\ref{thm:main}, sufficiency]
Let $P=\prod_iv_i/\prod_ic_i>1$ and $s=P^{1/N}>1$, so that $s^N=P$. Put $a_i=v_i/(c_is)>0$. Then
$\prod_ia_i=\prod_iv_i/(\prod_ic_i\cdot s^N)=1$. Define the prefix products
$x_i=\prod_{j=0}^{i}a_j>0$. For $i\ge1$, $x_i=a_i\,x_{i-1}=a_i\,x_{\mathrm{prev}(i)}$. For $i=0$,
$x_{\mathrm{prev}(0)}=x_{N-1}=\prod_ja_j=1$, so again $x_0=a_0=a_0\,x_{\mathrm{prev}(0)}$. Hence for
every $i$,
\[
\mathrm{net}_i=v_ix_{\mathrm{prev}(i)}-c_ia_ix_{\mathrm{prev}(i)}
=v_ix_{\mathrm{prev}(i)}\Bigl(1-\frac1s\Bigr)>0. \qedhere
\]
\end{proof}

\begin{remark}[Reading the construction]
Writing $\rho_i=x_i/x_{\mathrm{prev}(i)}$, partner $i$ benefits if and only if $\rho_i<v_i/c_i$.
Around the cycle the ratios must multiply to $\prod_i\rho_i=1$. A beneficial flow therefore exists
exactly when $1<\prod_iv_i/c_i$, which is the product condition. The construction splits the
surplus factor $P$ evenly in geometric terms: every link gets $\rho_i=(v_i/c_i)/s$. This is
\emph{one} feasible choice. We do not claim it is fair or efficient, or that partners would agree to
it. For $N=2$ the per-link condition $\rho_i<v_i/c_i$ is exactly Run-128's corridor
$c_X/v_X<y/x<v_Y/c_Y$.
\end{remark}

\begin{remark}[Which hypotheses are used]
Necessity does not use $v_i>0$, and it uses $N\ge2$ only to get a non-empty index set. The
equivalence also holds for $N=1$, where the cycle is a self-loop and the condition reads $c_0<v_0$.
It fails for $N=0$, where every flow is vacuously beneficial but $1<1$ is false. Sufficiency uses
$c_i,v_i>0$. The strict inequality is needed: the sealed module's control
$c=v=(1,1)$ shows that equality of the products is not enough.
\end{remark}

\paragraph{Example.} The sealed module's three-partner example has $c=(1,1,1)$ and $v=(4,1/2,1)$.
Partner~1 values its input at only half its unit cost, yet $\prod c=1<2=\prod v$. The uniform flow
leaves partner~1 with net benefit $-1/2$. The tuned flow $(1,9/20,2/5)$ gives net benefits
$(3/5,1/20,1/20)$. Our construction gives $s=2^{1/3}$ and flow
$x=(4/s,\;2/s^2,\;2/s^3)=(4/s,\;2/s^2,\;1)$. Each partner keeps the fraction
$1-2^{-1/3}\approx0.206$ of the value it receives.

\section{The \texorpdfstring{$N=2$}{N = 2} case: recovering the Reciprocity Corridor}
\label{sec:n2}
Instantiate the cycle with $N=2$, $c=(c_X,c_Y)$, $v=(v_X,v_Y)$ and $x=(x,y)$. Then
$\mathrm{prev}(0)=1$ and $\mathrm{prev}(1)=0$. The net benefits are Run-128's $u_X=v_Xy-c_Xx$ and
$u_Y=v_Yx-c_Yy$, and $\prod c<\prod v$ is Run-128's \texttt{ProductCondition}. The candidate file
contains a corollary, \texttt{two\_partner\_corridor}, that derives from \texttt{main} the
two-partner feasibility statement
\[
\bigl(\exists\,x,y>0:\ v_Xy-c_Xx>0\ \wedge\ v_Yx-c_Yy>0\bigr)\iff c_Xc_Y<v_Xv_Y .
\]
Through Run-128's Theorem~1 (benefit iff corridor) this is the content of its Theorem~2 (the
corridor is non-empty iff the product condition holds). The corollary is a helper in the candidate
file and not a sealed statement, so the verifier's label applies only to \texttt{main}.

\section{The Lean development}
\label{sec:proof}
The candidate file has five parts. \texttt{prev\_val} gives the explicit form of the predecessor by
case split and \texttt{Nat.add\_mod\_right}. \texttt{prev\_bijective} gets bijectivity from
injectivity on a finite type, and \texttt{prod\_prev} reindexes with
\texttt{Fintype.prod\_bijective}. \texttt{necessity} uses
\texttt{Finset.prod\_lt\_prod\_of\_nonempty}, \texttt{Finset.prod\_mul\_distrib} and
\texttt{lt\_of\_mul\_lt\_mul\_right}. \texttt{sufficiency} builds $s$ with \texttt{Real.rpow},
proves $s>1$ with \texttt{Real.one\_lt\_rpow} and $s^N=P$ with
\texttt{Real.rpow\_inv\_natCast\_pow}. It then defines the prefix products over
\texttt{Finset.range}, proves the cycle recursion $x_i=a_i\,x_{\mathrm{prev}(i)}$ and finishes with
\texttt{field\_simp} and \texttt{nlinarith}.

\begin{lstlisting}[caption={The complete candidate proof \texttt{proofs/P3.CyclicSymbiosis.lean}.}]
import Mathlib
import ViridisStatements.P3.CyclicSymbiosis

namespace Candidate

open ViridisStatements.P3.CyclicSymbiosis

/-- Explicit form of the predecessor: `0 ↦ N - 1`, `i ↦ i - 1` otherwise. -/
lemma prev_val {N : ℕ} (i : Fin N) :
    (prev i).val = if i.val = 0 then N - 1 else i.val - 1 := by
  have hN := i.isLt
  unfold prev
  split_ifs with h
  · simp only [h, zero_add]
    exact Nat.mod_eq_of_lt (by omega)
  · have : i.val + (N - 1) = (i.val - 1) + N := by omega
    simp only [this, Nat.add_mod_right]
    exact Nat.mod_eq_of_lt (by omega)

/-- The predecessor map is injective, hence a bijection of the finite cycle. -/
lemma prev_bijective {N : ℕ} : Function.Bijective (prev : Fin N → Fin N) := by
  rw [← Finite.injective_iff_bijective]
  intro a b hab
  have ha := prev_val a
  have hb := prev_val b
  have := congrArg Fin.val hab
  have := a.isLt
  have := b.isLt
  ext
  split_ifs at ha hb <;> omega

/-- Reindexing along the cycle: `∏ x (prev i) = ∏ x i`. -/
lemma prod_prev {N : ℕ} (x : Fin N → ℝ) : ∏ i, x (prev i) = ∏ i, x i :=
  Fintype.prod_bijective prev prev_bijective _ _ (fun _ => rfl)

/-- Necessity: a strictly beneficial flow forces `∏ c < ∏ v`. -/
lemma necessity (N : ℕ) (hN : 2 ≤ N) (c v : Fin N → ℝ) (hc : ∀ i, 0 < c i)
    (x : Fin N → ℝ) (hx : StrictlyBeneficialFlow c v x) : CycleProductCondition c v := by
  obtain ⟨hxpos, hnet⟩ := hx
  have hne : (Finset.univ : Finset (Fin N)).Nonempty := ⟨⟨0, by omega⟩, Finset.mem_univ _⟩
  have hlt : ∏ i, (c i * x i) < ∏ i, (v i * x (prev i)) := by
    apply Finset.prod_lt_prod_of_nonempty _ _ hne
    · intro i _; exact mul_pos (hc i) (hxpos i)
    · intro i _
      have := hnet i
      unfold net at this
      linarith
  rw [Finset.prod_mul_distrib, Finset.prod_mul_distrib, prod_prev] at hlt
  have hX : 0 < ∏ i, x i := Finset.prod_pos (fun i _ => hxpos i)
  exact lt_of_mul_lt_mul_right hlt hX.le

/-- Sufficiency: under `∏ c < ∏ v` an explicit flow benefits every partner. -/
lemma sufficiency (N : ℕ) (hN : 2 ≤ N) (c v : Fin N → ℝ) (hc : ∀ i, 0 < c i)
    (hv : ∀ i, 0 < v i) (hP : CycleProductCondition c v) :
    ∃ x : Fin N → ℝ, StrictlyBeneficialFlow c v x := by
  unfold CycleProductCondition at hP
  have hCpos : 0 < ∏ i, c i := Finset.prod_pos (fun i _ => hc i)
  have hVpos : 0 < ∏ i, v i := Finset.prod_pos (fun i _ => hv i)
  set P : ℝ := (∏ i, v i) / (∏ i, c i) with hPdef
  have hP1 : 1 < P := by rw [hPdef, one_lt_div hCpos]; exact hP
  set s : ℝ := P ^ ((N : ℝ)⁻¹) with hsdef
  have hs1 : 1 < s := Real.one_lt_rpow hP1 (by positivity)
  have hs0 : 0 < s := by linarith
  have hsN : s ^ N = P := Real.rpow_inv_natCast_pow (by linarith) (by omega)
  -- per-partner ratios, with product one
  set a : Fin N → ℝ := fun i => v i / (c i * s) with hadef
  have hapos : ∀ i, 0 < a i := fun i => div_pos (hv i) (mul_pos (hc i) hs0)
  have haprod : ∏ i, a i = 1 := by
    simp only [hadef, Finset.prod_div_distrib, Finset.prod_mul_distrib, Finset.prod_const,
      Finset.card_univ, Fintype.card_fin, hsN, hPdef]
    field_simp
  -- extend to ℕ and take prefix products
  set a' : ℕ → ℝ := fun j => if h : j < N then a ⟨j, h⟩ else 1 with ha'def
  have ha'pos : ∀ j, 0 < a' j := by
    intro j; simp only [ha'def]; split_ifs
    · exact hapos _
    · exact one_pos
  set y : ℕ → ℝ := fun k => ∏ j ∈ Finset.range (k + 1), a' j with hydef
  have hypos : ∀ k, 0 < y k := fun k => Finset.prod_pos (fun j _ => ha'pos j)
  set x : Fin N → ℝ := fun i => y i.val with hxdef
  -- the key recursion around the cycle
  have hrec : ∀ i, x i = a i * x (prev i) := by
    intro i
    have hpv := prev_val i
    have hi := i.isLt
    simp only [hxdef]
    split_ifs at hpv with h0
    · rw [hpv, h0]
      have hyN : y (N - 1) = 1 := by
        simp only [hydef]
        rw [show N - 1 + 1 = N by omega, ← haprod, ← Fin.prod_univ_eq_prod_range]
        apply Finset.prod_congr rfl
        intro j _
        simp [ha'def, j.isLt]
      rw [hyN, mul_one]
      simp only [hydef, zero_add, Finset.prod_range_one, ha'def]
      rw [dif_pos (by omega)]
      congr 1
      ext; simp [h0]
    · rw [hpv]
      simp only [hydef]
      rw [show i.val + 1 = (i.val - 1 + 1) + 1 by omega, Finset.prod_range_succ,
        show i.val - 1 + 1 = i.val by omega, mul_comm]
      congr 1
      simp only [ha'def, dif_pos hi]
  refine ⟨x, fun i => hypos _, fun i => ?_⟩
  unfold net
  rw [hrec i]
  have hX : 0 < x (prev i) := hypos _
  have hci := hc i
  have hvi := hv i
  have : c i * (a i * x (prev i)) = v i * x (prev i) / s := by
    simp only [hadef]; field_simp
  rw [this]
  have hvx : 0 < v i * x (prev i) := mul_pos hvi hX
  rw [sub_pos, div_lt_iff₀ hs0]
  nlinarith

theorem main : ViridisStatements.P3.CyclicSymbiosis.Statement := by
  intro N hN c v hc hv
  exact ⟨fun ⟨x, hx⟩ => necessity N hN c v hc x hx, sufficiency N hN c v hc hv⟩

/-- `N = 2` corollary: specialising `main` to the two-partner cycle gives Run-128's strict
feasibility test (Theorem 2): strictly positive flows `x, y` with both payoffs
`uX = vX*y - cX*x` and `uY = vY*x - cY*y` strictly positive exist iff
`ProductCondition vX cX vY cY`. -/
theorem two_partner_corridor (vX cX vY cY : ℝ) (hvX : 0 < vX) (hcX : 0 < cX)
    (hvY : 0 < vY) (hcY : 0 < cY) :
    (∃ x y : ℝ, 0 < x ∧ 0 < y ∧ 0 < vX * y - cX * x ∧ 0 < vY * x - cY * y) ↔
      ProductCondition vX cX vY cY := by
  have h0 : prev (0 : Fin 2) = 1 := by decide
  have h1 : prev (1 : Fin 2) = 0 := by decide
  have hc : ∀ i, 0 < (![cX, cY] : Fin 2 → ℝ) i := by
    intro i; fin_cases i <;> simp [hcX, hcY]
  have hv : ∀ i, 0 < (![vX, vY] : Fin 2 → ℝ) i := by
    intro i; fin_cases i <;> simp [hvX, hvY]
  have key := main 2 le_rfl ![cX, cY] ![vX, vY] hc hv
  have hcond : CycleProductCondition ![cX, cY] ![vX, vY] ↔ ProductCondition vX cX vY cY := by
    simp [CycleProductCondition, ProductCondition, Fin.prod_univ_two]
  rw [← hcond, ← key]
  constructor
  · rintro ⟨x, y, hx, hy, hX, hY⟩
    refine ⟨![x, y], ?_, ?_⟩
    · intro i; fin_cases i <;> simp [hx, hy]
    · intro i
      fin_cases i
      · simpa [net, h0] using hX
      · simpa [net, h1] using hY
  · rintro ⟨f, hpos, hnet⟩
    refine ⟨f 0, f 1, hpos 0, hpos 1, ?_, ?_⟩
    · simpa [net, h0] using hnet 0
    · simpa [net, h1] using hnet 1

end Candidate
\end{lstlisting}

\section{Results}
\label{sec:results}
We ran the program's checker on the candidate file. It runs the same stages as the verifier:
denylist, compile, statement match, axioms and kernel replay. The output was
\texttt{PASSED (preview)}, and the reported axioms were \texttt{propext},
\texttt{Classical.choice} and \texttt{Quot.sound}. The first version of the proof passed on its
first attempt. We then added the $N=2$ corollary, and the file passed again. Compiling it with
\texttt{lean} gave no warnings. No automated prover was used. The certification page that precedes
this paper gives the verifier's own judgement.

With this result all four statements planned for P3 (loop mass balance, cascaded recovery, the
recovery-adjusted work floor and cyclic symbiosis) have candidate proofs that passed the preview
check.

\section{What this settles and what it does not}
\label{sec:limits}
\paragraph{Settled.} Assuming the verifier accepts the file, the product condition is exactly the
feasibility condition for strictly beneficial flows on a cycle, for every $N\ge2$ and all positive
real coefficients.

\paragraph{Not settled.}
\begin{itemize}
\item \textbf{No real coefficients.} Nothing here estimates any real cost $c_i$ or value $v_i$.
Whether a real loop satisfies the condition is an empirical question, and the answer stays
assumed.
\item \textbf{Model assumptions.} Payoffs are linear and separable, there is one period, there are
no fixed costs, no capacity limits and no externalities, and coefficients are known. Fixed
costs, transport and storage logistics, capacity limits or quality loss along the loop could all
change the answer.
\item \textbf{Contracts and bargaining.} The theorem says a beneficial flow \emph{exists}. It does
not say which flow partners would agree on, how the surplus would be split, or how an agreement
would be enforced. Run-128 studies a Nash-bargaining split for $N=2$. Nothing like that is
proved here for $N>2$.
\item \textbf{Only simple cycles.} General exchange networks (several suppliers, branching flows,
several materials) are not covered. Generalising to a strongly connected exchange graph, where
we conjecture the condition becomes a cycle-mean or spectral-radius condition, is open.
\end{itemize}

\section{Reproducibility}
\sloppy
The proof file \texttt{proofs/P3.CyclicSymbiosis.lean} imports only \texttt{Mathlib} and the
sealed module \texttt{ViridisStatements.P3.CyclicSymbiosis}. It compiles under toolchain
\texttt{leanprover/lean4:v4.28.0} and the Mathlib commit below.
\begin{lstlisting}[language={},frame=none]
sealed module sha256:
  a6825700fb4c21d12ba330c739d321dd105da3b21c394338cda11ef5a67667c4
Mathlib commit:
  8f9d9cff6bd728b17a24e163c9402775d9e6a365
viridis-lean-check proofs/P3.CyclicSymbiosis.lean P3.CyclicSymbiosis
\end{lstlisting}

\section{Changelog}
\begin{itemize}
\item 2026-11-05 (P3-T001): first version. Proposed P3.LoopMassBalance, P3.CascadedRecovery,
P3.RecoveryWorkFloor and P3.CyclicSymbiosis as statement modules, with controls, probe results
and rationale.
\item 2026-11-15 (P3-T002): P3.LoopMassBalance sealed (D-011). Narrowed the paper to a note on
this statement with a full Lean proof and a human-readable proof.
\item 2026-11-16 (P3-T003): P3.CascadedRecovery sealed (D-012). Note on the cascade statement with
a full Lean proof.
\item 2026-11-17 (P3-T004): P3.RecoveryWorkFloor sealed (D-013). Note on the recovery-adjusted
work floor with a full Lean proof built on Run-163.
\item 2026-11-18 (P3-T005): P3.CyclicSymbiosis sealed (D-014). This version is a note on the
$N$-partner product condition, with a full Lean proof, an explicit beneficial flow and a Lean
corollary that recovers the Run-128 two-partner test. Added the Run-128 Zenodo reference. Earlier
results are summarised in the background section; their proofs stay in the earlier run packages.
\end{itemize}

\bibliographystyle{plainnat}
\bibliography{refs}
\end{document}
